\section{تمرین عملی}

آنچه را باید بسازند توضیح دهید و مستندات مرجع را پیوند بزنید:

\url{https://example.com/docs}

دامنه‌ی کار را روشن بگویید: چه چیزی لازم است، چه چیزی خارج از محدوده است، و
کدام داده‌ها می‌توانند فقط در حافظه بمانند.

\subsection{مراحل}

\begin{enumerate}
    \item مستندات مرجع را بخوانید.
    \item پیش از نوشتن کد، رویکردتان را توضیح دهید.
    \item پیاده‌سازی کنید.
    \item معماری نهایی را تحلیل کنید.
\end{enumerate}

\subsection{یک پنجره‌ی کد}

\begin{neocode}[neocodetitlefont=\pixelfont]{config.yaml}
listen-on: 80

paths:
  "/":
    type: redirect
    target: "http://localhost:8080/index.html"
  "/dir/":
    type: servedir
    target: "/home/user/movies/"
\end{neocode}

\lr{\texttt{neocode}} متن نوار عنوان را به عنوان آرگومان اجباری می‌گیرد و هر
کلید \lr{\texttt{listings}} را در آرگومان اختیاری، مثلاً
\lr{\texttt{[listing options=\{language=Python\}]}}. داخل بلوک،
\lr{\texttt{@\ldots@}} به لاتک بازمی‌گردد، پس ماکروهایی مانند
\lr{\texttt{\textbackslash hwnumber}} همچنان باز می‌شوند.

\subsection{فهرست تیک‌دار}

\begin{neochecklist}
    \citem موردی که انجام شده، با \lr{\texttt{\textbackslash citem}}
    \item موردی که انجام نشده، با \lr{\texttt{\textbackslash item}}
    \item موردی دیگر که پیش از تحویل باید انجام شود
\end{neochecklist}
